\title{Parameter-Efficient Orthogonal Finetuning via Factorization}

\begin{abstract}
Large foundation models have become widespread, but training them from the ground up remains prohibitively costly. Consequently, efficiently adapting these powerful models to specific downstream tasks has gained importance. In this work, we investigate a well-founded finetuning strategy—Orthogonal Finetuning (OFT)—for task adaptation. Although OFT exhibits strong generalization capabilities, it still requires a considerable number of trainable parameters due to the large dimensionality inherent in orthogonal matrices. To tackle this challenge, we first analyze OFT through the lens of information transmission and then pinpoint several crucial criteria that promote improved parameter efficiency. Drawing inspiration from the Cooley-Tukey fast Fourier transform algorithm, which facilitates efficient information transmission, we introduce an efficient orthogonal parameterization employing butterfly structures. Incorporating this parameterization into OFT, we develop a novel parameter-efficient finetuning technique termed Orthogonal Butterfly (BOFT). BOFT generalizes OFT by encompassing it as a special case, thereby establishing a broader orthogonal finetuning framework. Finally, we perform comprehensive empirical evaluations adapting large vision transformers, language models, and text-to-image diffusion models to diverse downstream tasks in both vision and language domains.
\end{abstract}

\section{Introduction}

Recent models like ChatGPT~\cite{brown2020language,chen2021evaluating} and Stable Diffusion~\cite{rombach2022high} showcase the impressive generalization capabilities of large foundation models. The swift advancements in their performance have been accompanied by a significant escalation in parameter counts (e.g., GPT-3 contains about 175B parameters). Consequently, training foundation models from scratch has become increasingly difficult. Progress in the field thus hinges on the ability to adapt these models efficiently without full retraining. In other words, we need methods for effectively adapting existing foundation models to downstream tasks. Efficient task adaptation primarily falls into three categories: \emph{model finetuning}~\cite{hulora2022,zaken2022bitfit,guo2020parameter,raffel2020exploring,qiu2023controlling,zhang2022adaptive,chavan2023one}, where the architecture remains fixed and only a subset of parameters is updated; \emph{adapter tuning}~\cite{rebuffi2017learning,houlsby2019parameter,pfeiffer2020adapterfusion,lin2020exploring,he2021towards}, which introduces additional trainable parameters into the model; and \emph{prompt tuning}~\cite{li2021prefix,lester2021power}, where trainable prefix tokens are prepended to inputs. Among these, model finetuning stands out as a straightforward yet effective method that importantly does not increase inference latency.

The core idea behind model finetuning is to maintain similarity between the pretrained model and the finetuned model according to certain metrics, thereby preserving the knowledge acquired during pretraining. For example, many existing finetuning techniques employ a small learning rate. This heuristic method helps to minimize the difference between the pretrained and finetuned models. Considering the structured properties of weight matrices, a more principled strategy aims to retain the relational information within these matrices, specifically the pairwise angles between neurons. This concept gives rise to a new finetuning framework called Orthogonal Finetuning~(OFT)~\cite{qiu2023controlling}, where neurons within the same layer are transformed using a common orthogonal matrix, ensuring that the pairwise angles between neurons remain invariant throughout the finetuning process. Although OFT has shown promising results in terms of generalizability and convergence when finetuning text-to-image diffusion models~\cite{qiu2023controlling}, the number of trainable parameters can be quite large due to the high dimensionality of the orthogonal matrices. To mitigate this, OFT employs a block-diagonal structure to reduce parameter count. However, this gain in parameter efficiency comes with drawbacks—the orthogonal matrix features a fixed sparsity pattern and the orthogonal transformation is applied independently across different blocks. This block-diagonal sparsity pattern, while parameter-efficient, may introduce certain undesirable inductive biases (e.g., limiting expressiveness and inability to approximate classical linear transformations), and more importantly, determining an optimal sparsity pattern remains an open challenge.

The crucial aspect of solving this problem lies in generating a dense orthogonal matrix while maintaining parameter efficiency. Although this might initially appear impractical because a $d$-dimensional dense orthogonal matrix demands $\mathcal{O}(d^2)$ parameters, we introduce a novel method that constructs a dense orthogonal matrix by composing multiple factorized sparse matrices. This strategy is motivated by the insight that the number of trainable parameters can be decreased by exchanging computation time for memory usage. Since the orthogonal matrix is represented as the product of sparse matrices, the multiplication must be executed repeatedly during each training iteration. To better understand this matrix factorization, we reinterpret the creation of a dense orthogonal matrix as an information transmission problem. Within this framework, generating a dense orthogonal matrix through multiplication of sparse matrices can be viewed as transmitting information over a grid-structured graph. This perspective allows us to devise multiple ways to conduct sparse matrix factorization that restrict the number of trainable parameters while retaining sufficient expressiveness to produce dense matrices.

To achieve parameter efficiency within our information transmission framework, we draw inspiration from the butterfly structures found in the Cooley-Tukey fast Fourier transform algorithm~\cite{cooley1965algorithm}, which facilitates efficient information exchange among nodes~\cite{klasing1994broadcasting}. In particular, the butterfly graph intrinsic to the Cooley-Tukey algorithm naturally induces a sparse matrix factorization method known as butterfly factorization. Utilizing butterfly factorization, we can generate a $d$-dimensional dense matrix as a product of $\mathcal{O}(\log d)$ sparse matrices, each containing $\mathcal{O}(d)$ non-zero elements. By ensuring that each sparse matrix is orthogonal, we obtain an orthogonal parameterization that is highly efficient, requiring only $\mathcal{O}(d \log d)$ parameters, which markedly reduces the original parameter count. Building on this efficient orthogonal parameterization, we introduce a novel parameter-efficient finetuning technique termed Orthogonal Butterfly (BOFT). BOFT generalizes OFT as a special case, offering a comprehensive orthogonal finetuning framework. A common trait shared by both the block-diagonal and butterfly structures is sparsity; both structures exploit data sparsity in orthogonal matrices to decrease the effective number of trainable parameters. It is also insightful to compare our approach with the low-rank structure utilized in LoRA~\cite{hulora2022}; both low-rank and sparse matrices represent prominent families of structured matrices~\cite{candes2011robust} employed to achieve parameter efficiency.

In contrast to the block-diagonal structure used by OFT to balance expressiveness and regularity, BOFT employs the butterfly structure to create a smoother interpolation between matrices from the full orthogonal group (expressivity) and identity matrices (regularity). This approach allows us to identify a smaller hypothesis space within the orthogonal group suitable for downstream tasks. Given the extensive use of the butterfly structure in various fast linear transforms, such as the discrete Fourier and discrete cosine transforms, we contend that our structured strategy for parameter efficiency introduces an implicit inductive bias that enhances generalizability and mitigates overfitting. Our intuition stems from the fact that the butterfly structure can readily reconstruct many classical linear transforms, whereas the block-diagonal structure in OFT cannot recover any. Our main contributions are summarized as follows:

\begin{itemize}[leftmargin=*,nosep]
    \item We explore the challenge of parameter efficiency in orthogonal finetuning through a novel information transmission framework. This framework reframes the task of designing a parameter-efficient dense orthogonal matrix as an information transmission problem within a grid-structured graph.
    \item Drawing inspiration from the butterfly structures in the Cooley-Tukey algorithm, we introduce Orthogonal Butterfly, a parameter-efficient orthogonal finetuning approach, formulated within the information transmission framework.
    \item We offer several theoretical insights explaining why BOFT can substantially reduce the number of trainable parameters while maintaining strong expressivity and training stability. Owing to matrix factorization, BOFT also features an intriguing weight interpolation (see Figure~\ref{fig:boft_interpolation}).
    \item For the first time, we extend orthogonal finetuning to a variety of tasks beyond controllable text-to-image generation~\cite{qiu2023controlling}, demonstrating its significant potential as a universal model finetuning technique. Specifically, we apply BOFT across numerous adaptation scenarios spanning computer vision and natural language processing. BOFT surpasses current state-of-the-art methods by a notable margin, confirming its superior parameter efficiency and generalization capability.
\end{itemize}

\section{Related Work}

\textbf{Parameter-efficient finetuning (PEFT)}. With foundation models (\emph{e.g.}, \cite{brown2020language,radford2021learning,rombach2022high,kirillov2023segment}) growing larger and more capable, there has been significant interest in finetuning these models for downstream tasks in a parameter-efficient manner~\cite{treviso2023efficient,ding2023parameter,lialin2023scaling}. Among various PEFT techniques~\cite{houlsby2019parameter,aghajanyan2020intrinsic,hulora2022,edalati2022krona,wang2022adamix,gheini2021cross,zaken2022bitfit,guo2020parameter,sung2021training,ansell2022composable,lester2021power,li2021prefix,vu2022spot,he2021towards,mao2021unipelt,karimi2021compacter,liu2022few,sung2022lst,chen2022parameter,jia2022visual,chen2022adaptformer,zhang2022neural,jie2023fact,lian2022scaling,luo2023towards}, reparameterization-based approaches~\cite{aghajanyan2020intrinsic,hulora2022,edalati2022krona} are the most pertinent to our study. LoRA~\cite{hulora2022} updates pretrained weight matrices by adding the product of two low-rank matrices, yielding strong performance on natural language tasks. Since LoRA fixes the rank of all added matrices, several works \cite{zhang2022adaptive,valipour2022dylora,zhang2023increlora} propose dynamically adjusting the rank layer-wise to better allocate parameter budgets. Due to its straightforwardness, this low-rank weight reparameterization approach has become widely adopted \cite{dettmers2023qlora,zi2023delta,chavan2023one}. Drawing inspiration from how hyperspherical energy influences generalization~\cite{liu2018learning,liu2021orthogonal}, \cite{qiu2023controlling} introduces orthogonal finetuning (OFT), an alternative yet effective method to finetune text-to-image diffusion models. In particular, OFT learns an orthogonal matrix that transforms neurons within the same layer, achieving better generalization and more stable training compared to LoRA. Although OFT performs well, it typically involves more trainable parameters than LoRA. Hence, improving the parameter efficiency of OFT is a valuable objective. Additionally, it remains unclear if OFT can be extended beyond controlling text-to-image diffusion models~\cite{qiu2023controlling} to a broader range of adaptation tasks. BOFT enhances OFT’s parameter efficiency through butterfly factorization, enabling us to showcase the efficacy of orthogonal finetuning for general adaptation tasks.

\textbf{Butterfly structures}. The radix-2 Cooley-Tukey algorithm recursively decomposes an $N$-point discrete Fourier transform into two $\frac{N}{2}$-point discrete Fourier transforms, producing a butterfly structure that can be expressed as a product of multiple sparse matrices (known as a butterfly matrix). Butterfly matrices have been employed to parameterize orthogonal matrices to avoid pivoting during Gaussian elimination and improve computational efficiency~\cite{parker1995random}, to stabilize recurrent neural network training~\cite{jing2017tunable}, and in kernel approximation~\cite{munkhoeva2018quadrature}. Works such as \cite{dao2019learning,dao2019kaleidoscope} learn fast linear transforms using butterfly parameterizations. Other studies \cite{chen2022pixelated,dao2022monarch} utilize butterfly matrices to facilitate efficient neural network training. Furthermore, butterfly structures prove beneficial in fast matrix-vector multiplication~\cite{michielssen1996multilevel,o2010algorithm}, data-sparse matrix approximation~\cite{li2015butterfly}, and network transmission~\cite{klasing1994broadcasting,li2003linear,rajasingh2016transmission}. Distinct from prior work, our focus is on leveraging butterfly structures to improve the parameter efficiency of OFT.

\section{An Information Transmission View on Orthogonal Finetuning}

We begin with some foundational concepts of OFT. To finetune the pretrained weight matrix, OFT re-expresses the new weight matrix as the product of a learnable orthogonal matrix and the fixed pretrained weight matrix. In contrast to LoRA, which updates weights by adding a low-rank matrix, OFT applies a multiplicative orthogonal matrix for the updates. For parameter efficiency, LoRA utilizes a low-rank design, whereas the original OFT adopts a block-diagonal orthogonal structure~\cite{qiu2023controlling} (the smaller the diagonal blocks, the more parameter-efficient the OFT). An intuitive comparison is illustrated in Figure~\ref{fig:comp_lora}. The rationale behind using an orthogonal transformation to finetune the weight matrix is to maintain the pairwise angles between neurons~\cite{liu2018learning,liu2021orthogonal,qiu2023controlling}, thus largely preserving the semantic information acquired during pretraining. Specifically, OFT optimizes an orthogonal matrix $\bm{R}\in\mathbb{R}^{d\times d}$ for a pretrained linear layer $\bm{W}^0\in\mathbb{R}^{d\times n}$, modifying the forward computation from $\bm{z}=(\bm{W}^0)^\top\bm{x}$ to $\bm{z}=(\bm{R}\bm{W}^0)^\top\bm{x}$, where $\bm{x}\in\mathbb{R}^{d}$ and $\bm{z}\in\mathbb{R}^{n}$ denote the input and output vectors, respectively. To realize zero initialization, OFT sets $\bm{R}$ initially as the identity matrix. To preserve orthogonality of $\bm{R}$ throughout finetuning, we adopt the Cayley parameterization as in \cite{qiu2023controlling,liu2021orthogonal}, i.e., $\bm{R}=(\bm{I}+\bm{Q})(\bm{I}-\bm{Q})^{-1}$, where $\bm{Q}$ is a skew-symmetric matrix satisfying $\bm{Q}=-\bm{Q}^\top$. For parameter efficiency, the block-diagonal structure is introduced by representing the orthogonal matrix $\bm{R}$ as $\text{diag}(\bm{R}_1,\bm{R}_2,\cdots,\bm{R}_r)$, where each $\bm{R}_i\in\mathcal{R}^{b\times b}, \forall i$ is a small orthogonal matrix and $br=d$. The parameter efficiency afforded by this block-diagonal design comes with a trade-off — it assumes that the dimensions of a neuron (i.e., a column vector in $\bm{W}^0$) are split into $r$ groups, and that different groups are transformed independently with distinct orthogonal matrices. While empirically effective, this block-diagonal structure lacks a meaningful basis for dividing neuron dimensions into $r$ groups by their indices, suggesting that dense orthogonal matrices are preferable. This raises a natural question: \emph{Is it possible to construct a dense orthogonal matrix that maintains parameter efficiency?}

To tackle this issue, we suggest constructing a dense orthogonal matrix by expressing it as a product of several sparse orthogonal matrices. To offer a cohesive and intuitive framework for examining the sparsity patterns in orthogonal matrix factorization, we interpret the task of generating a dense orthogonal matrix in OFT as an information transmission problem. More precisely, creating a dense matrix $\bm{R}\in\mathbb{R}^{d\times d}$ through the product of $m$ square matrices, $\thickmuskip=2mu \medmuskip=2mu \bm{R}=\bm{B}_m\bm{B}_{m-1}\cdots\bm{B}_1$, can be viewed as conveying information across a grid composed of $d\times (m+1)$ nodes, as depicted in Figure~\ref{fig:info_trans}. This perspective is motivated by the recognition that a $d$-dimensional dense square matrix represents dense connectivity between $d$ nodes and another set of $d$ nodes. Regarding the matrix $\bm{R}$, a zero element $R_{ij}$ signifies that information from the $j$-th node cannot reach the $i$-th node, whereas a non-zero $R_{ij}$ indicates that such transmission is possible. Consequently, expressing the dense matrix $\bm{R}$ as a product $\bm{B}_m\bm{B}_{m-1}\cdots\bm{B}_1$ can equivalently be seen as executing a sequential information exchange guided by the graphs associated with each $\bm{B}_i, \forall i$. Information propagates starting with $\bm{B}_1$ and concluding with $\bm{B}_m$. For instance, in Figure~\ref{fig:info_trans}, we examine the factorization $\bm{R}=\bm{B}_5\bm{B}_4\bm{B}_3\bm{B}_2\bm{B}_1$, with its sparsity patterns and corresponding graph displayed. This graph results from unrolling the matrix multiplication. Within this induced graph, each matrix $\bm{B}_i$ represents the connectivity from nodes at level $i$ to nodes at level $i+1$. Specifically, the $(j_1,j_2)$-th entry of $\bm{B}_i$ indicates whether there is a directed edge from the $j_2$-th node at level $i$ to the $j_1$-th node at level $i+1$ (with zero meaning no edge). For the product $\bm{B}_5\bm{B}_4\bm{B}_3\bm{B}_2\bm{B}_1$ to form a dense matrix, every node at the first level must be capable of transmitting information to all nodes at the sixth level. Considering only the product $\bm{R}=\bm{B}_4\bm{B}_3\bm{B}_2\bm{B}_1$, which links source nodes at the first level to receiver nodes at the fifth level, reveals that information from node 1 cannot be sent to node 3. Thus, the sparsity pattern of $\bm{B}_4\bm{B}_3\bm{B}_2\bm{B}_1$ contains a zero at position $(3,1)$. Similar correspondences between the induced graph and sparsity patterns hold when considering $\bm{R}=\bm{B}_3\bm{B}_2\bm{B}_1$ and $\bm{R}=\bm{B}_2\bm{B}_1$. In general, to obtain a dense matrix $\bm{R}\in\mathbb{R}^{d\times d}$, the $m$ factor matrices $\bm{B}_m, \ldots, \bm{B}_1$ must correspond to directed edges arranged on a $d\times (m+1)$ grid, where each directed edge connects nodes only between adjacent levels (i.e., columns), ensuring that information from every node in the first level can be transmitted to every node in the $(m+1)$-th level.

The information transmission perspective aids in better comprehending the sparsity pattern of factorization matrices in OFT. Figure~\ref{fig:blk_diag} illustrates the block-diagonal configuration of $\bm{R}$ in the original OFT. Although this block-diagonal form decreases the number of trainable parameters, it cannot produce a dense matrix $\bm{R}$. Our objective is to construct a dense orthogonal matrix by combining $m$ sparse orthogonal matrices, while minimizing the number of effective trainable parameters. From the information transmission standpoint, the key requirements for this goal are \emph{(i)} \textbf{dense connectivity}: each node at the first level must have at least one path to every node at the last level, and \emph{(ii)} \textbf{minimum free edges}: the total edge count should be as low as possible under the orthogonality constraint. It is important to recognize that orthogonality imposes a subtle constraint on the edges between adjacent levels. For instance, for each matrix $\bm{B}_i$ to be full-rank (a necessary condition for orthogonality), there must be $d$ edges to establish a bijection between all nodes at the $i$-th level and all nodes at the $(i+1)$-th level, ensuring the number of edges between these levels is at least $d$ (\emph{e.g.}, 4 in the example shown in Figure~\ref{fig:info_trans}). These $d$ edges are essential for orthogonality and should not be included in the edge count since these elements are not trainable (\emph{e.g.}, in a $d \times d$ orthogonal matrix with $d$ non-zero entries, these entries can only be $\pm 1$). Given that orthogonal matrices require fewer parameters than full matrices, the orthogonality condition introduces additional dependencies among edges. For example, in a $2 \times 2$ orthogonal matrix, having a zero at the $(1,1)$ position implies another zero at the $(2,2)$ position (\emph{i.e.}, the absence of one edge may necessitate the absence of another). Consequently, for each feasible edge connection set, orthogonality may sometimes add or eliminate edges. Visualizing the non-zero pattern of the composed orthogonal matrix makes the information transmission framework especially useful in OFT, since only the non-zero trainable elements of $\bm{R}$ are relevant, and their exact values do not matter.

A straightforward dense connection linking two levels requires $\mathcal{O}(d^2)$ edges (i.e., a single dense orthogonal matrix), resulting in $d^2 - d$ edges (for instance, when $d=4$, this equates to 12 edges). Figure~\ref{fig:info_trans} presents an example of a valid matrix factorization that uses a total of $10$ edges, which is actually fewer than what is needed for one dense orthogonal matrix. This framework allows us to analyze the parameter efficiency of OFT through a graphical lens, making it easy to identify feasible factorizations. Our approach is inspired by a notable topology from the Cooley-Tukey algorithm, known as butterfly graphs~\cite{cooley1965algorithm}, which can efficiently create dense connections between $d$ source nodes and $d$ receiver nodes with only $\mathcal{O}(d \log d)$ edges. For example, the topology shown in Figure~\ref{fig:info_trans} employs $10$ edges to achieve dense connectivity, whereas the butterfly network requires just $8$ edges. We now discuss how the butterfly structure enhances parameter efficiency.

\section{Orthogonal Parameterization by Butterfly Factorization}

Originally devised for the Cooley-Tukey algorithm to accelerate the Fourier transform, the butterfly structure reflects the principle that a local change in the frequency domain induces a global change in the spatial domain—paralleling our information transmission problem, where every node at the initial level can convey information to all nodes at the final level. The butterfly structure has also become a popular computer network topology~\cite{solihin2015fundamentals,li2011linear} utilized for efficient information exchange. Assuming $k \geq 2$ is a power of $2$, we define the \emph{butterfly factor} $\bm{B}^F(k)$ as

\begin{equation}\label{eq:but_factor}
\begin{aligned}
    \bm{B}^F(k)=\begin{bmatrix}
    \text{diag}(\bm{d}_1) & \text{diag}(\bm{d}_2) \\
    \text{diag}(\bm{d}_3) & \text{diag}(\bm{d}_4)
    \end{bmatrix} \in \mathbb{R}^{k \times k},
\end{aligned}
\end{equation}

where each vector $\bm{d}_i \in \mathbb{R}^{\frac{k}{2}}$ for $i = 1,\dots,4$. For $d = 2^N$, we recursively define the $d$-dimensional \emph{butterfly matrix} $\bm{B}(d) \in \mathbb{R}^{d \times d}$ as

\begin{equation}
\begin{aligned}
    \bm{B}(d) = \tilde{\bm{B}}(d,d) \cdot \begin{bmatrix}
    \bm{B}_1\left(\frac{d}{2}\right) & \bm{0} \\
    \bm{0} & \bm{B}_2\left(\frac{d}{2}\right)
    \end{bmatrix} = \tilde{\bm{B}}(d,d) \tilde{\bm{B}}(d,\frac{d}{2}) \cdots \tilde{\bm{B}}(d,2),
\end{aligned}
\end{equation}

where $\bm{B}_1(\frac{d}{2})$ and $\bm{B}_2(\frac{d}{2})$ represent two butterfly matrices each of dimension $\frac{d}{2}$. Next, we introduce the \emph{butterfly component} defined as $\thickmuskip=2mu \medmuskip=2mu \tilde{\bm{B}}(d,k)=\text{diag}(\bm{B}^F_1(k),\cdots,\bm{B}^F_{d/k}(k))$, which is a block-diagonal matrix of size $d\times d$ with blocks of size $k$, where each $\bm{B}^F_i(k), \forall i$, corresponds to butterfly factors given by Equation~\ref{eq:but_factor}. We are now prepared to use the butterfly matrix to represent an orthogonal matrix. To do this, it suffices to ensure that each multiplicative factor $\tilde{\bm{B}}(d,k), \forall k$ within the butterfly matrix $\bm{B}(d)$ is orthogonal. Consider first the block-diagonal matrix $\tilde{\bm{B}}(d,2)$ with block size 2; it can be easily made orthogonal by parameterizing each block using the Cayley transform (or a $2$-dimensional rotation), as done in \cite{liu2021orthogonal,qiu2023controlling}. The sparsity pattern of every butterfly component corresponds to a permutation of the non-zero structure of $\tilde{\bm{B}}(d,2)$, hence all butterfly components can be parameterized straightforwardly as orthogonal matrices. This yields an efficient orthogonal matrix parameterization constructed from multiple $2 \times 2$ orthogonal blocks. Extending this idea following \cite{chen2022pixelated}, we define a block butterfly component $\tilde{\bm{B}}^b(d,k)$ where each scalar entry in $\bm{d}_i, \forall i$, is replaced by a $b \times b$ matrix. To ensure the block butterfly component $\tilde{\bm{B}}^b(d,2)$ is orthogonal, each $2b \times 2b$ block is parameterized to be orthogonal. The sparsity pattern of the other block butterfly components $\tilde{\bm{B}}^b(d,k), k>2$, is a block-wise permutation of that of $\tilde{\bm{B}}^b(d,2)$ and can thus be similarly made orthogonal. Putting all parts together, the forward pass in BOFT is

\begin{equation}
    \begin{aligned}\nonumber
    \bm{z}=\big(\bm{R}(m,b)\cdot\bm{W}^0\big)^\top\bm{x},~~~~\text{s.t.}~\bigg{\{}\bm{R}(m,b)=\prod_{i=1}^m \tilde{\bm{B}}^b_{(d,i)}~~\&~~\underbrace{\big(\tilde{\bm{B}}^b_{(d,j)}\big)^\top\tilde{\bm{B}}^b_{(d,j)}=\tilde{\bm{B}}^b_{(d,j)}\big(\tilde{\bm{B}}^b_{(d,j)}\big)^\top\!=\bm{I}_d}_{\forall j\in [1, m]}\bigg{\}},
    \end{aligned}
\end{equation}

where, for simplicity, we represent $\tilde{\bm{B}}^b(d,2^{m - i+1})$ as $\tilde{\bm{B}}^b_{(d,i)}$, and $\bm{I}_d$ stands for the $d \times d$ identity matrix. The orthogonal matrix $\bm{R}(m,b)\in\mathbb{R}^{d\times d}$ is formed by the product of multiple orthogonal butterfly components. For ease of notation, we refer to BOFT with $\bm{R}(m,\frac{b}{2})$ as BOFT($m$,$b$), where $b\geq 2$. When $m=1$, BOFT($1$,$b$) simplifies to the block-diagonal OFT~\cite{qiu2023controlling} with block size $b$. Similarly, BOFT($1$,$d$) corresponds to the original OFT~\cite{qiu2023controlling}, featuring an unconstrained full orthogonal matrix. BOFT($\log \frac{2d}{b}$,$b$) employs the block butterfly matrix $\bm{B}^{\frac{b}{2}}(d)$ as $\bm{R}$, resulting in a dense orthogonal matrix $\bm{R}$. Generally, BOFT($m$,$b$) involves $\frac{1}{2}(b-1)dm$ trainable parameters for fine-tuning a linear layer of size $d\times n$. When using the butterfly matrix with $m=\log d$ and $b=2$, BOFT requires $\mathcal{O}(d\log d)$ parameters. In comparison, the original OFT with a full dense orthogonal matrix uses $\mathcal{O}(d^2)$ parameters, while the block-diagonal OFT with $r$ blocks uses $\mathcal{O}(bd)$. Thus, the original OFT must use block size $b=d$ to produce a dense orthogonal matrix, whereas BOFT can achieve this with any $b$.

\textbf{Identity initialization for BOFT}. Fine-tuning approaches typically begin from the exact pretrained model to ensure the fine-tuned model remains close to the original. For instance, LoRA initializes the low-rank weights to zero. In BOFT, all butterfly components are initialized as identity matrices (i.e., the skew-symmetric matrix is set to zero in the Cayley transform).

\textbf{Multiplicative Dropout for BOFT}. LoRA~\cite{hulora2022} incorporates a Dropout layer on the low-rank weight updates to mitigate overfitting. While standard Dropout~\cite{srivastava2014dropout} is naturally applicable to LoRA, it is not directly compatible with BOFT due to the multiplicative nature of its weight updates. To overcome this, we introduce a multiplicative Dropout mechanism tailored for BOFT. Given that the orthogonal matrix $\bm{R}(m,b)$ consists of $m$ orthogonal butterfly components, which can be rearranged into $2b \times 2b$ block-diagonal orthogonal matrices, this multiplicative Dropout randomly selects $p_1$ percent of the butterfly components and $p_2$ percent of the diagonal blocks within each component, substituting those blocks with identity matrices.

\section{Intriguing Insights and Discussions}

\textbf{Expressivity of BOFT}. The combination of the butterfly structure and permutations can exactly replicate numerous classic fast linear transforms~\cite{dao2019learning,dao2019kaleidoscope} (such as the fast Fourier transform and Hadamard transform). However, the extent to which our orthogonal butterfly matrix can approximate a general orthogonal matrix is still unknown. To investigate this, we perform a simulation to approximate a random dense orthogonal matrix~\cite{anderson1987generation} of size $512 \times 512$. The results shown in Figure~\ref{fig:para_eff} are averaged over 10 different random seeds. The y-axis represents the approximation error, while the x-axis indicates the number of effective trainable parameters. Each curve of the same color corresponds to BOFT with a fixed block size, and the leftmost point represents the error of BOFT($1$,$b$) (i.e., the original block-diagonal OFT with block size $b$). Generally, BOFT achieves superior parameter efficiency compared to OFT. For instance, BOFT($9$,$2$) demonstrates better expressiveness than BOFT($1$,$16$) while using significantly fewer parameters. BOFT configurations with smaller $b$ and larger $m$ tend to be more parameter-efficient. For example, BOFT($6$,$4$) employs far fewer parameters but attains a similar approximation error as BOFT($2$,$16$). Overall, the butterfly matrix corresponds to a more structured subset of the orthogonal group (relative to the block-diagonal structure), which makes BOFT provably more expressive than OFT with an identical block size.

\begin{theorem}[Expressivity of BOFT]\label{thm:exp_boft}
    BOFT exhibits greater expressivity than OFT when the block size is the same. To approximate all orthogonal matrices of size $d$ using butterfly matrices, one can multiply matrices in the form $\bm{B}_{d-1,1}(d)\bm{B}^\top_{d-1,2}(d)\cdots\bm{B}_{1,1}(d)\bm{B}^\top_{1,2}(d)$, where each $\bm{B}_{i,j}(d)$ is a butterfly matrix for all $i$ and $j$.
\end{theorem}

Theorem~\ref{thm:exp_boft} motivates a straightforward generalization of BOFT—the final orthogonal matrix can be expressed as $\thickmuskip=2mu \medmuskip=2mu \bm{R}^G(m_1,b_1,m_2,b_2,l)=\bm{R}_{l,1}(m_1,b_1)\bm{R}_{l,2}^T(m_2,b_2)\cdots\bm{R}_{1,1}(m_1,b_1)\bm{R}_{1,2}^T(m_2,b_2)$, where $\bm{R}_{i,j}^T(m,b)$ denotes the orthogonal matrix employed in BOFT. When setting $m_1 = m_2 = \log d$, $b_1 = b_2 = 2$, and $l = d - 1$, the matrix $\bm{R}^G(m_1,b_1,m_2,b_2,l)$ can represent the entire orthogonal group. This matrix composition is also known as the kaleidoscope hierarchy~\cite{dao2019kaleidoscope}. Nevertheless, it is important to note that increased expressivity does not always translate to improved performance during finetuning. For instance, full finetuning, despite its universal expressiveness, often produces unsatisfactory results. Achieving the right balance between expressivity and regularity is crucial for the generalizability of model finetuning. BOFT expands the finetuning parameter space by incorporating structural priors, enabling a better compromise between expressivity and regularity.

\textbf{Spectral properties}. Orthogonal finetuning typically provides superior spectral characteristics compared to LoRA because it exactly maintains the spectral norm of the pretrained weight matrix $\bm{W}^0$. This can be understood through singular value decomposition: $\bm{W}^0=\bm{U}\bm{\Sigma}\bm{V}^\top$, where $\bm{U}$ and $\bm{V}$ are orthogonal matrices and $\bm{\Sigma}$ is a diagonal matrix of singular values. Both OFT and BOFT apply an orthogonal matrix $\bm{R}$ on the left side, resulting in the finetuned weights $\bm{R}\bm{U}\bm{\Sigma}\bm{V}^\top$, which does not alter the largest singular value (i.e., the spectral norm of $\bm{W}^0$). This preservation has been demonstrated to significantly enhance training stability and generalization~\cite{miyato2018spectral,yoshida2017spectral}. Additional intriguing mathematical properties are discussed in Appendix~\ref{app:sn_property}.

\textbf{Orthogonal finetuning as learning bilinear similarity}. BOFT can be interpreted as learning the bilinear similarity $\bm{w}^0_i \bm{R} \bm{x}$, where $\bm{w}^0_i$ denotes the $i$-th neuron (i.e., column vector) of the weight matrix $\bm{W}^0$. Essentially, BOFT involves learning the bilinear similarity matrix $\bm{R}$ under a strong constraint (i.e., $\bm{R}$ must be orthogonal), which inherently relates to distance metric learning~\cite{xing2002distance} and bilinear forms~\cite{roman2005advanced}.

\textbf{Inductive bias and generalization in BOFT}. Since $\bm{R}(m,b)$ in BOFT typically corresponds to a structured subset of the orthogonal group that restricts the hypothesis space, BOFT naturally introduces an inductive bias. We contend that the structured inductive bias introduced by butterfly factorization is advantageous for generalization, as it embodies a common structured pattern found in many classical linear transforms~\cite{dao2019kaleidoscope}, including the discrete Fourier transform, discrete sine/cosine transform, and Hadamard transform. Furthermore, the sparse matrix factorization inherent in BOFT may also contribute implicit inductive biases~\cite{gunasekar2017implicit,li2018algorithmic,liu2019neural}.

\textbf{Comparison to butterfly-based sparse training}. Several studies~\cite{dao2019learning,dao2019kaleidoscope,chen2022pixelated,dao2022monarch} have explored sparse training utilizing butterfly parameterization. These works generally concentrate on directly reparameterizing weight matrices with the butterfly structure and training neural networks from the beginning. In \cite{dao2022monarch}, the focus is on finetuning pretrained weights by first projecting them onto a variant of butterfly matrices, followed by optimizing the projected components for downstream tasks. BOFT introduces a fundamentally different finetuning approach that modifies the weights using layer-shared weight matrices.

\input{rebuttal-results}

\section{Concluding Remarks and Limitations}

This paper presents Orthogonal Butterfly, a versatile and parameter-efficient finetuning approach for foundation models founded on the butterfly architecture. The crucial idea to enhance parameter efficiency is to represent a dense orthogonal matrix as a product of several sparse orthogonal matrices. To facilitate finding suitable matrix factorizations, we develop a graphical information transmission framework. Within this framework, we identify that the butterfly structure effectively fulfills our criteria for sparse orthogonal matrix factorization. We validate the practical benefits of BOFT through finetuning experiments on large language models, large vision models, and text-to-image generative models. Our results also demonstrate the superiority of BOFT as a broadly applicable finetuning technique.

While BOFT exhibits strong empirical performance, it is not without limitations. Because the final orthogonal matrix in BOFT is composed as the product of multiple orthogonal matrices, the training runtime overhead is slightly higher compared to OFT. Finding ways to reduce BOFT's training runtime overhead remains an open challenge. On the upside, after the finetuning phase, the orthogonal matrices learned by BOFT can be directly integrated into the pretrained model without incurring extra inference latency. Additionally, it remains unclear whether the butterfly network is the optimal topology for information transmission. Our information transmission framework opens avenues to draw insights from another field—computer networking—where the efficiency of network topologies for information transfer is extensively studied. We anticipate that alternative, more efficient network structures could be utilized for constructing dense orthogonal matrices.

\end{document}